\documentclass[10pt,a4paper]{amsart}
\usepackage[margin=19mm]{geometry}
\usepackage{amsmath,amssymb,mathtools}
\usepackage[scr=rsfs,cal=euler]{mathalfa}
\usepackage{xfrac}
\usepackage[unicode,hidelinks,bookmarks=false]{hyperref}
\newcommand{\ie}{i.e.\ }
\newcommand{\cf}{cf.\ }
\newcommand{\resp}{respectively\ }
\newcommand{\Subsection}{\subsection}
\providecommand{\nobreakdash}{\nobreak\hbox{-}}
\setlength{\emergencystretch}{2em}
\multlinegap0pt
\usepackage{amssymb}
\newtheorem{lemma}{Lemma}
\title{Every quasiperfect number has at least eight distinct prime factors}
\author{Akira Toyohara, Ye Tao, Siqiong Yao}
\date{Source: arXiv:2608.02066v1 (CC BY 4.0)}
\begin{document}
\maketitle

\noindent\emph{Source and licence.} Every quasiperfect number has at least eight distinct prime factors. Version arXiv:2608.02066v1 (CC BY 4.0), distributed by arXiv under the Creative Commons Attribution 4.0 International licence, http://creativecommons.org/licenses/by/4.0/, as recorded in the arXiv licence metadata for this exact version and shown on the abstract page https://arxiv.org/abs/2608.02066v1. Full licence notice: \texttt{licenses/arxiv-2608.02066-CC-BY-4.0.txt}.\par\smallskip

\noindent\textit{Editorial scope.} Counted unit: the article's lemma lem:disc (source0.tex lines 160--165). The article's complete original proof of it is delivered with it (source0.tex lines 167--172, one \texttt{proof} environment of 6 lines). No mathematics is written, repaired or re-derived: the statement and the proof are the article's own lines, in the article's own notation, and no retained line is substituted or re-flowed.  Retained uncounted as the source-local context the proof needs: lines 97--98.  Nothing was omitted from any retained span, so the package carries no omission record.  No retained line contains a citation, so the wrapper carries no bibliography and no bibliography entry.  No retained line contains a figure environment, an image-inclusion command or drawing code, so no image is delivered and the package declares no asset dependency.  Editorial formatting only, in addition: an AMS \texttt{amsart} preamble that restates the article's own declarations and macros used by the retained lines, namely the environments \texttt{lemma} on separate counters, together with the standard packages those lines need.  The retained environments print in this document as Lemma 1 in order of appearance, whereas the article prints the same items with the numbering of its own full document.  The complete native source of the article as distributed in the arXiv e-print for this exact version is bundled unchanged as \texttt{sources/206664/source0.tex}.
\begin{equation} A\,\sigma(p^2)\,\sigma(q^2) = 2Np^2q^2 + 1. \label{eq:leaf}
\end{equation} Here $N, A$ are the known constants of the leaf, and the only unknowns left are the two primes $p, q$.

\begin{lemma}[discriminant criterion]\label{lem:disc} Let $N, A, p$ be positive integers, and write $S = A(p^2 + p + 1)$, $T = 2Np^2$. If a positive integer $q$ satisfies the leaf equation \eqref{eq:leaf}, then
\[ \Delta(p) := S^2 + 4(T - S)(S - 1)
\] is a perfect square, and moreover
\begin{equation} q = \frac{S + \sqrt{\Delta(p)}}{2(T - S)}. \label{eq:qformula}
\end{equation}
\end{lemma}

\begin{proof} With $\sigma(q^2) = q^2 + q + 1$, substitute into \eqref{eq:leaf} and collect in $q$:
\[ S(q^2 + q + 1) = Tq^2 + 1 \iff (T - S)\,q^2 - S\,q - (S - 1) = 0.
\] This is a quadratic in $q$. Multiplying both sides by $4(T - S)$ and completing the square gives
\[ \bigl(2(T - S)q - S\bigr)^2 = S^2 + 4(T - S)(S - 1) = \Delta(p).
\] The left-hand side is a perfect square, so $\Delta(p)$ is a perfect square. Taking the positive root gives formula \eqref{eq:qformula}.
\end{proof}

\end{document}
